% SOURCE_AUTHORITY: sga5_fr_workpass.tex, lines 15042--15066 (LF-normalized unit).
% SOURCE_SHA256: 791F4EFFC5E02832D5D77ED03518C8156D6F07E4C8238B03545DB93D883FBB28
\begin{statement}{Lema 3}
Sea $X$ un esquema de tipo finito sobre $\mathbb F_p$, y sea $g:X\to e$ su morfismo estructural.
\begin{enumerate}[label=\alph*)]
\item Consideremos un $\Omega$-haz constructible $F$ sobre $X$, un subesquema cerrado $Y$ de $X$ y el abierto complementario $U=X-Y$. Si el teorema es cierto para dos de los pares $(X,F)$, $(Y,F|Y)$, $(U,F|U)$, es cierto para el tercero.
\item Consideremos una sucesión exacta
\[
0\longrightarrow F'\longrightarrow F\longrightarrow F''\longrightarrow0
\]
de $\Omega$-haces constructibles sobre $X$; si el teorema es cierto para dos de los pares $(X,F')$, $(X,F)$, $(X,F'')$, es cierto para el tercero.
\end{enumerate}
\end{statement}

Solo utilizaremos el enunciado a), pero damos también b) por simetría. Las demostraciones de a) y b) son análogas; demostremos, por ejemplo, a). La sucesión exacta
\[
\cdots\longrightarrow R^i_!g_U(F|U)\longrightarrow R^i_!g(F)
\longrightarrow R^i_!g_Y(F|Y)\longrightarrow R^{i+1}_!g_U(F|U)\longrightarrow\cdots
\]
y la proposición 1a) (n\textsuperscript{o}~1) dan
\[
\prod_i\bigl(L_{R^i_!g(F)}\bigr)^{(-1)^i}
=\prod_i\bigl(L_{R^i_!g_U(F|U)}\,L_{R^i_!g_Y(F|Y)}\bigr)^{(-1)^i},
\]
y la proposición 1b) (n\textsuperscript{o}~1) permite entonces concluir.

Los razonamientos de los lemas 1, 2 y 3 pueden resumirse diciendo que, según la proposición 1 del n\textsuperscript{o}~1, la función $L_F$ tiene propiedades formales análogas a las de la cohomología con soportes propios.
